%%% Zhejiang University of Technology 
%%% !Mode\dots "TeX:UTF-8"
%%% author : liuzheng712@gmail.com(刘正)

\pagenumbering{arabic}          % 页码数字
\setcounter{page}{1}            % 页码起始数字
\renewcommand{\sfdefault}{phv}
\definecolor{dkgreen}{rgb}{0,0.6,0}
\definecolor{gray}{rgb}{0.5,0.5,0.5}
\newcommand{\tab}{\hspace{3em}}

\definecolor{matlab-string}{RGB}{160, 32, 240}
\definecolor{matlab-keyword}{RGB}{0, 0, 255}
\definecolor{matlab-comment}{RGB}{34, 139, 34}
\definecolor{xml-constant}{RGB}{0, 0, 0}

\newcommand{\ms}[1]{\textcolor{matlab-string}{#1}}
\newcommand{\mk}[1]{\textcolor{matlab-keyword}{#1}}
\newcommand{\mc}[1]{\textcolor{matlab-comment}{#1}}

\newcommand{\xt}[1]{\textcolor{matlab-keyword}{<#1>}}
\newcommand{\xa}[1]{\textcolor{matlab-string}{#1}}
\newcommand{\xs}[1]{\textcolor{matlab-string}{"#1"}}
\newcommand{\xc}[1]{\textcolor{matlab-comment}{<!-- #1 -->}}
\renewcommand{\xi}[1]{\textcolor{xml-constant}{#1}}

\newenvironment{code}{\alltt\small}{\endalltt}
\newenvironment{boxedcode}{\begin{boxedminipage}{\linewidth}\begin{code}}{\end{code}\end{boxedminipage}}

% 行距 1.25
\renewcommand{\baselinestretch}{1.25}
\newcommand{\tablecell}[2]{\begin{tabular}{@{}#1@{}}#2\end{tabular}}   % 表格内换行，用法\tablecell{c}{balabalba\\balabalba}
\graphicspath{{figures/}}  		% 图路径 

%%%%%%字体部分
\newfontfamily\fangsong{Adobe Fangsong Std}     % 仿宋
\newfontfamily\hei{Adobe Heiti Std}             % 黑体
\newfontfamily\kai{Adobe Kaiti Std}             % 楷体
\newfontfamily\song{Adobe Song Std}             % 宋体
\newfontfamily\Bsong{华文中宋}             % 粗宋体
\newfontfamily\monaco{Monaco}
\newfontfamily\times{Times New Roman}
\newcommand{\shijian}{\number\year~年~\number\month~月~\number\day~日}
\renewcommand{\textipa}[1]{\fontencoding{T3}\selectfont #1 \times}
\font\yinbiao="DejaVu Serif" at 11pt      % 修复重音符号bug
%\renewcommand{\t}[1]{\zhongyinfuhao \t{#1} \times}
%\renewcommand{\b}[1]{\zhongyinfuhao \b{#1} \times}
%\renewcommand{\c}[1]{\zhongyinfuhao \c{#1} \times}
%\renewcommand{\.}[1]{\zhongyinfuhao \.{#1} \times}
%\renewcommand{\r}[1]{\zhongyinfuhao \r{#1} \times}
%\renewcommand{\t}[1]{\zhongyinfuhao #1 \times}
%\renewcommand{\t}[1]{\zhongyinfuhao #1 \times}

%%%%%%%字号%%%%%%%
%%%%%%单倍%%%%%%%%
\newcommand{\yihao}{\fontsize{26pt}{26pt}\selectfont}       % 一号, 单倍行距
\newcommand{\xiaoyi}{\fontsize{24pt}{24pt}\selectfont}      % 小一, 单倍行距
\newcommand{\erhao}{\fontsize{22pt}{22pt}\selectfont}       % 二号, 单倍行距
\newcommand{\xiaoer}{\fontsize{18pt}{18pt}\selectfont}      % 小二, 单倍行距
\newcommand{\sanhao}{\fontsize{16pt}{16pt}\selectfont}      % 三号, 单倍行距
\newcommand{\xiaosan}{\fontsize{15pt}{15pt}\selectfont}     % 小三, 单倍行距
\newcommand{\sihao}{\fontsize{14pt}{14pt}\selectfont}       % 四号, 单倍行距
\newcommand{\xiaosi}{\fontsize{12pt}{12pt}\selectfont}      % 小四, 单倍行距
\newcommand{\wuhao}{\fontsize{10.5pt}{10.5pt}\selectfont}   % 五号, 单倍行距
\newcommand{\xiaowu}{\fontsize{9pt}{9pt}\selectfont}        % 小五, 单倍行距

%%%%%%1.5倍%%%%%%%
\newcommand{\Yihao}{\fontsize{26pt}{39pt}\selectfont}       % 一号, 1.5倍行距
\newcommand{\Xiaoyi}{\fontsize{24pt}{36pt}\selectfont}      % 小一, 1.5倍行距
\newcommand{\Erhao}{\fontsize{22pt}{33pt}\selectfont}       % 二号, 1.5倍行距
\newcommand{\Xiaoer}{\fontsize{18pt}{27pt}\selectfont}      % 小二, 1.5倍行距
\newcommand{\Sanhao}{\fontsize{16pt}{24pt}\selectfont}      % 三号, 1.5倍行距
\newcommand{\Xiaosan}{\fontsize{15pt}{22.5pt}\selectfont}   % 小三, 1.5倍行距
\newcommand{\Sihao}{\fontsize{14pt}{21pt}\selectfont}       % 四号, 1.5倍行距
\newcommand{\Xiaosi}{\fontsize{12pt}{18pt}\selectfont}      % 小四, 1.5倍行距
\newcommand{\XIaosi}{\fontsize{12pt}{15pt}\selectfont}      % 小四, 1.25倍行距
\newcommand{\Wuhao}{\fontsize{10.5pt}{15.75pt}\selectfont}  % 五号, 1.5倍行距
\newcommand{\Xiaowu}{\fontsize{9pt}{13.5pt}\selectfont}     % 小五, 1.5倍行距

%%%%%%%%%% Table, Figure and Equation %%%%%%%%%%%%%%%%%
\renewcommand{\tablename}{\song\wuhao 表}                                     % 插表题头
\renewcommand{\figurename}{\song\wuhao 图}                                    % 插图题头
\renewcommand{\thefigure}{\arabic{chapter}-\arabic{figure}}       % 使图编号为 7-1 的格式 %\protect{~}
%\renewcommand{\thesubfigure}{\alph{subfigure})}                   % 使子图编号为 a) 的格式
%\renewcommand{\thesubtable}{(\alph{subtable})}                    % 使子表编号为 (a) 的格式
\renewcommand{\thetable}{\arabic{chapter}-\arabic{table}}         % 使表编号为 7-1 的格式
\renewcommand{\theequation}{\arabic{chapter}-\arabic{equation}}   % 使公式编号为 7-1 的格式
%%%%%%%%%%%%章节定义%%%%%%%%%%
\renewcommand{\contentsname}{\song\xiaoer 目\quad录} 
\renewcommand{\listfigurename}{\song\xiaoer 图\quad列} 
\renewcommand{\listtablename}{\song\xiaoer 表\quad列} 

%%%%%%%%%% Chapter and Section %%%%%%%%%%%%%
\setcounter{secnumdepth}{4}
\setlength{\parindent}{2em}
\renewcommand{\chaptername}{第\xCJKnumber{\thechapter}章}
\titleformat{\chapter}{\centering\erhao\hei\bfseries}{\hei\chaptername}{2em}{}
\titlespacing{\chapter}{0pt}{0.1\baselineskip}{0.8\baselineskip}

\titleformat{\section}{\sihao\hei\bfseries}{\hei\thesection}{1em}{}
\titlespacing{\section}{0pt}{\baselineskip}{0.3\baselineskip}

\titleformat{\subsection}{\sihao\hei\bfseries}{\hei\thesubsection}{1em}{}
\titlespacing{\subsection}{0pt}{0.1\baselineskip}{0.3\baselineskip}

\titleformat{\subsubsection}{\sihao\hei}{\thesubsubsection}{1em}{}
\titlespacing{\subsubsection}{0pt}{0.05\baselineskip}{0.1\baselineskip}

\newcommand{\cankao}[1]{$^{\mbox{\protect \scriptsize \cite{#1}}}$}% 修改引用文件样式

%\let\cleardoublepage\relax % delete the black page

%%%%%%%%%%%%%%%%%%%%%目录设定%%%%%%%%%%%%%%%%%%%%%%%%%%%%
%\titlecontents{chapter}[0em]{\vspace{1\baselineskip}\Xiaosi\times}{\Xiaosi\song 第\xCJKnumber{\thecontentslabel}章\quad}{}{\titlerule*[0.7pc]{.}\contentspage}
\titlecontents{chapter}[0em]{\Xiaosi\times}{\Xiaosi\song 第\xCJKnumber{\thecontentslabel}章\quad}{}{\titlerule*[0.7pc]{.}\contentspage}
\titlecontents{section}[2em]{\Xiaosi\times}{\Xiaosi\song\thecontentslabel\quad}{}{\titlerule*[0.7pc]{.}\contentspage}
\titlecontents{section}[2em]{\Xiaosi\times}{\Xiaosi\song\thecontentslabel\quad}{}{\titlerule*[0.7pc]{.}\contentspage}
\titlecontents{subsection}[4em]{\Xiaosi\times}{\Xiaosi\song\thecontentslabel\quad}{}{\titlerule*[0.7pc]{.}\contentspage}
%%%%%%%%%%%%%%%%%%%%%自定义目录形式%%%%%%%%%%%%%%%%%%%%%%%%%%%%%%%%%
%%%%%%%%%%%%%%%%%%%%%%%%自定义目录格式%%%%%%%%%%%%%%%%%%%
%1.定义形式如下4种
%\def\mlA#1#2#3{\par\noindent\hspace*{1 true cm}\parbox[b]{12.3 true cm}{\protect\hspace*{-1 true cm}{#1}\quad{#2}~\dotfill~}\hbox to1.0 true cm{~(\hss#3\hss)}\vspace{2mm}}
%/def/mlB#1#2#3{/par/noindent/hspace*{1 true cm}/parbox[b]{12.3 true cm}{/protect/hspace*{-1 true cm}{#1}~/dotfill~/hbox{#2}}/hbox to1.0 true cm{~(/hss#3/hss)}/vspace{2mm}}
%/def/mlC#1#2{/par/noindent/hspace*{1 true cm}/parbox[b]{12.3 true cm}{/protect/hspace*{-1 true cm}/quad{#1}~/dotfill~}/hbox to1.0 true cm{~(/hss#2/hss)}/vspace{2mm}}
%/def/mlD#1#2{/par/noindent/hspace*{1 true cm}/parbox[b]{12.3 true cm} {/protect/hspace*{-1 true cm}/quad{#1}~/hfill~}/hbox to1.0 true cm{~(/hss#2/hss)}/vspace{2mm}}
%2.使用方法如下：
%\mlA{title}{author}{page no.}
%效果： title author .................. (page no.)
%\mlB{title}{author}{page no.}
%效果：title .......................... author (page no.)
%\mlC{title}{page no.}
%效果： title ......................... (page no.)
%\mlD{title}{page no.}
%效果： title (page no.)
%%%%%%%%%%%%%%%%%%%%%%%%%%%%%%%%%%%%%%%%%%%%%%%%%%%%%%%%%%%%




%%%%%%%%%%%%%%%%%%%%%%%% Code %%%%%%%%%%%%%%%%%%%%%%%%%
\lstset{
        language=matlab,            % 设定默认语言为MATLAB
        keywords={break,case,catch,continue,else,elseif,end,for,function,
        global,if,otherwise,persistent,return,switch,try,while}, %设定关键词列表
        keywordstyle=\color{blue},  % 关键词为蓝色
        commentstyle=\color{dkgreen},       % 注释为绿色
        stringstyle=\color{red},    % 字符串为红色
        basicstyle=\wuhao\monaco, % 基本字体的字号
        breaklines=true,            % 自动将长代码行换行排版
        breakatwhitespace=true,     % 断行只在空格处
        extendedchars=false,        % 解决代码跨页时，章节标题页眉等汉字不显示问题
        showspaces=false,           % 不显示空格
        showstringspaces=true,      % 字符串中显示空格
        showtabs=false,             % 不显示TAB键
        tabsize=4,                  % TAB被当成4个空格
        frame=l,                    % 显示边框
        numbers=left,               % 显示行号
        numberstyle=\tiny,          % 行号字体为tiny
        numbersep=9pt,              % 行号垂直位置
        numberstyle=\normalsize,  % 行号字体的字号
        stepnumber=1,               % 行号显示的步长
        keywordstyle=\color{blue}\bfseries, % 特殊代码高亮蓝色加粗
        backgroundcolor=\color{white},      % 背景色 需要 \usepackage{color}
        escapeinside={/*@}{@*/}     % 添加注释，暂时离开
        }
\renewcommand{\lstlistingname}{CODE}
\lstloadlanguages{% Check Dokumentation for further languages ...
        MATLAB
}
\pagestyle{fancy}
\fancypagestyle{plain}

\def\yemeishow
{
\fancyhead{}                        % 清空页眉
\fancyhead[C]{\xiaowu\song\yemei}
\fancyfoot[C]{\xiaowu\song\thepage}
}
\def\yemeiclean
{
	\fancyhead{}                        % 清空页眉
	\renewcommand{\headrulewidth}{0pt}       %把页眉线的宽度设为零，即去掉页眉线	
}
%%%%%%%%%%%%%%%%%%%%%%%%%%%%%%%%%%%%%%%%%%%%%%%%%%%%%%%%%%%%%%%%%%%%%%%%%%%
%   首页
%%%%%%%%%%%%%%%%%%%%%%%%%%%%%%%%%%%%%%%%%%%%%%%%%%%%%%%%%%%%%%%%%%%%%%%%%%%
\def\makeZJUTthesisFirst
{
	\begin{center}
	\sanhao\hei
	\begin{figure}[h]
	\centerline{\includegraphics[height=3.2cm]{figures/zjutlogo.jpg}}
	\end{figure}
	
	\xiaoyi\hei 本科毕业设计论文\\[2.5cm]
	\end{center}
	\begin{center}
	
	\begin{quote}
	\hei\sanhao
	设计（论文）
	\\[0.5cm]
	题~~~~~~目\uline{\makebox[10cm]{\timu}}
	\\[1.5cm]
	指导教师
	\\
	姓~~~~~~名\uline{\makebox[10cm]{\laoshi}}
	\\[1.5cm]
	学~~~~~~生
	\\
	姓~~~~~~名\uline{\makebox[10cm]{\xuesheng}}
	\\[1.5cm]
	学~~~~~~生
	\\
	学~~~~~~号\uline{\makebox[10cm]{\xuehao}}
	\\[1.5cm]
	\xiaosan
	\begin{center}
	院~~~~~~系\uline{\makebox[6cm]{\xueyuan}}\\[1cm]
	专~~~~~~业\uline{\makebox[6cm]{\zhuanye}}\\[1cm]
	班~~~~~~级\uline{\makebox[6cm]{\banji}}\\[1cm]
	
	\end{center}
	\end{quote}
	\end{center}
	\thispagestyle{empty}
	\newpage
	\pagestyle{fancy}
\fancyhead{}                        % 清空页眉
\fancyhead[C]{\xiaowu\song\yemei}
\fancyfoot[C]{\xiaowu\song\thepage}
\fancypagestyle{plain}
}
%%%%%%%%%%%%%%%%%%%%%%%%%%%%%%%%%%%%%%%%%%%%%%%%%%%%%%%%%%%%%%%%%%%%%%%%%%%
%   文献综述的格式
%%%%%%%%%%%%%%%%%%%%%%%%%%%%%%%%%%%%%%%%%%%%%%%%%%%%%%%%%%%%%%%%%%%%%%%%%%%
\def\makeZJUTthesisWXZS
{

	\setcounter{chapter}{0}			% 章节起始
	\yemeiclean
	\pagenumbering{arabic}          % 页码数字
	\setcounter{page}{1}            % 页码起始数字
	\let\cleardoublepage\relax
	\begin{center}
	\sanhao\hei~
	\\[1cm]
	\xiaoyi\hei 浙江工业大学理学院
	
	毕业论文（设计）文献综述\\[4.5cm]
	\end{center}
	\begin{center}
	
	\begin{quote}
	\song\Sanhao
	\Bsong学生姓名：\uline{\song\makebox[3cm]{\xuesheng}}
	\Bsong学~~号：\song\uline{\makebox[4.6cm]{\xuehao}}
	
	\Bsong专~~~~~~业：\song{\uline{\makebox[10cm]{\zhuanye}}}
	
	\Bsong设计（论文）
	
	题~~~~~~目：\uline{\makebox[10cm]{\song\timu}}
	
	\Bsong{指导教师：}\uline{\makebox[10cm]{\song\laoshi}}
	
	\end{quote}
	\end{center}
	~\\[6cm]
	\begin{center}
	\song\xiaoer \shijian
	\end{center}
	\thispagestyle{empty}
	\newpage
	
	\renewcommand{\chaptername}{\xCJKnumber{\thechapter}、}
	\titleformat{\chapter}{\sihao\hei\bfseries}{\hei\chaptername}{0em}{}
	\titlespacing{\chapter}{0pt}{\baselineskip}{0pt}
	
	\titlespacing{\section}{0pt}{\baselineskip}{0.1\baselineskip}
	
	\pagenumbering{arabic}          % 页码数字
	\setcounter{page}{1}            % 页码起始数字
	\newpage
}

%%%%%%%%%%%%%%%%%%%%%%%%%%%%%%%%%%%%%%%%%%%%%%%%%%%%%%%%%%%%%%%%%%%%%%%%%%%
%   开题报告的格式
%%%%%%%%%%%%%%%%%%%%%%%%%%%%%%%%%%%%%%%%%%%%%%%%%%%%%%%%%%%%%%%%%%%%%%%%%%%
\def\makeZJUTthesisKTBG
{
	\yemeiclean
	\setcounter{chapter}{0}			% 章节起始
	
	\pagenumbering{arabic}          % 页码数字
	\setcounter{page}{1}            % 页码起始数字
	\let\cleardoublepage\relax
	\begin{center}
	\sanhao\hei~
	\\[1cm]
	\xiaoyi\hei 浙江工业大学理学院
	
	毕业论文（设计）开题报告\\[4.5cm]
	\end{center}
	\begin{center}
	
	\begin{quote}
	\song\Sanhao
	\Bsong学生姓名：\uline{\song\makebox[3cm]{\xuesheng}}
	\Bsong学~~号：\song\uline{\makebox[4.6cm]{\xuehao}}
	
	\Bsong专~~~~~~业：\song{\uline{\makebox[10cm]{\zhuanye}}}
	
	\Bsong设计（论文）
	
	题~~~~~~目：\uline{\makebox[10cm]{\song\timu}}
	
	\Bsong{指导教师：}\uline{\makebox[10cm]{\song\laoshi}}
	
	\end{quote}
	\end{center}
	~\\[6cm]
	\begin{center}
	\song\xiaoer \shijian
	\end{center}
	\thispagestyle{empty}
	\newpage
	
	\renewcommand{\chaptername}{\xCJKnumber{\thechapter}、}
	\titleformat{\chapter}{\sihao\hei\bfseries}{\hei\chaptername}{0em}{}
	\titlespacing{\chapter}{0pt}{\baselineskip}{0pt}
	
	\titlespacing{\section}{0pt}{\baselineskip}{0.1\baselineskip}
	
	\pagenumbering{arabic}          % 页码数字
	\setcounter{page}{1}            % 页码起始数字
	\newpage
}
\yemeishow
